% -------------------------------------------------------
\section{The Recursive Spawning Protocol}
\label{sec:recursive-spawning-protocol}
\index{The Recursive Spawning Protocol}

\subsection{Overview and Design Motivation}

The Recursive Spawning Protocol (RSP) is the operational core of the BX3 Framework's accountable multi-agent architecture. It defines the precise sequence of interactions between the Purpose Layer, Bounds Engine, and Fact Layer that occurs each time one agent authorizes the creation of another. RSP converts the abstract requirement of immutable parent chaining into a deterministic, enforceable, and auditable lifecycle event.

The need for RSP arises from a fundamental tension in hierarchical multi-agent systems: task decomposition requires agents to spawn children that operate with delegated authority, but every spawned child that lacks a traceable authorization path to a human principal is an agent operating in an state of \emph{autonomous drift} (Section~\ref{sec:autonomous-drift}). Recursive spawning amplifies this risk geometrically, not linearly: a single drifted parent with branching factor $b$ produces up to $b^d$ drifted descendants at depth $d$. The standard industry response — depth limits — does not address this amplification because drift can occur at any depth, and a drifted agent at depth $d=1$ is as dangerous as one at depth $d=10$ (Section~\ref{sec:drift-problem}).

RSP resolves this tension by making every spawn event a formally structured, three-layer operation in which:

\begin{itemize}
    \item The Purpose Layer validates that the spawn request is authorized and consistent with the parent's declared mandate.
    \item The Bounds Engine verifies that the proposed depth is within configured limits and that the child falls within convergence criteria.
    \item The Fact Layer creates an immutable parent pointer, records the spawn event to the forensic ledger, and activates the child with a sealed parent reference.
    \item The child is activated with a structurally immutable parent pointer, preventing any subsequent re-parenting or parent-severing events.
    \item The full chain is cryptographically verifiable as terminating at a human anchor.
\end{itemize}

The result is that recursive spawning, which in conventional architectures is a provenance gap, becomes the most precisely audited event in the system.

% -------------------------------------------------------
\subsection{Protocol Specification}

RSP is formally defined as a four-phase atomic protocol executed across the three BX3 layers. The protocol is atomic in the database sense: all four phases succeed together, or the spawn event is rejected in entirety and no partial record is inscribed. There is no spawn event that reaches the forensic ledger without complete Purpose Layer authorization.

\bigskip
\noindent\textbf{Phase 1: Purpose Layer Validation.}

The parent agent $p$ submits a \texttt{SpawnRequest} to the Purpose Layer. The request contains:

\begin{itemize}
    \item \textbf{Child Identifier:} A proposed unique identifier $c$ for the child agent.
    \item \textbf{Mandate:} A formal description of the task, scope, and authorization boundaries proposed for $c$.
    \item \textbf{Parent Attestation:} A cryptographic signature binding the request to $p$'s own authorized mandate, demonstrating that $p$ has the authority to delegate the specified scope.
    \item \textbf{Proposed Depth:} The depth $d_c = d_p + 1$ at which $c$ will operate.
\end{itemize}

The Purpose Layer performs three validation checks before proceeding:

\begin{enumerate}
    \item \textbf{Authorization Validity:} The parent $p$ must hold an active authorization that includes the right to delegate the proposed scope. The parent's own mandate is retrieved from its sealed Worksheet, and the proposed delegation scope must be a proper subset of what $p$ itself is authorized to act upon. This prevents escalation-of-authority drift: a parent cannot delegate what it does not itself hold.
    \item \textbf{Mandate Coherence:} The proposed mandate for $c$ must be interpretable as a consistent sub-goal of $p$'s parent mandate. If the Purpose Layer detects that the proposed mandate is orthogonal to or in contradiction with $p$'s mandate, the spawn is rejected. This prevents goal mutation from propagating through the spawn chain.
    \item \textbf{Termination Check:} The Purpose Layer verifies that the task decomposition represented by this spawn event has a defined termination condition. If the task assigned to $c$ does not have a specified convergence criterion, the spawn is flagged for human review before authorization is granted.
\end{enumerate}

If all three checks pass, the Purpose Layer issues a \texttt{SpawnAuthorization} token signed with the parent's Purpose Layer key. This token is the sole input to Phase 2.

\bigskip
\noindent\textbf{Phase 2: Bounds Engine Depth and Convergence Verification.}

Upon receiving a valid \texttt{SpawnAuthorization} token, the Bounds Engine independently verifies two conditions:

\begin{enumerate}
    \item \textbf{Maximum Depth Enforcement:} The proposed depth $d_c$ is compared against the configured system parameter $\lambda_{\text{max\_depth}}$. If $d_c > \lambda_{\text{max\_depth}}$, the spawn is refused. The Bounds Engine maintains this parameter independently of the Purpose Layer to ensure depth limits are enforced as a deterministic constraint, not a policy judgment. The system also enforces a minimum depth clearance: if $d_c \leq 0$, the request is malformed and rejected.
    \item \textbf{Convergence Criterion Verification:} The Bounds Engine verifies that the child $c$'s mandate includes at least one formally specified convergence criterion — a condition that, when satisfied, triggers the child's orderly termination and return of results to $p$. This prevents the spawn of children that have no defined exit condition, which is the structural cause of orphan drift (Section~\ref{sec:drift-failure-modes}). A child without a convergence criterion is a potential infinite descendant chain.
\end{enumerate}

If either condition fails, the Bounds Engine returns a \texttt{SpawnRefusal} to the Purpose Layer, which logs the refusal event to the Fact Layer and returns an error to $p$. No immutable record of the refused child is created; the refusal is recorded as a denied authorization attempt.

\bigskip
\noindent\textbf{Phase 3: Fact Layer Immutable Pointer Creation and Forensic Ledger Recording.}

Upon successful completion of Phases 1 and 2, the Fact Layer executes the atomic write that creates the spawn event:

\begin{enumerate}
    \item \textbf{Parent Pointer Assignment:} The Fact Layer records the binary relation $ParentPointer(p, c)$ as a permanent, single-assignment fact. This relation is immutable after inscription — no operation in the protocol permits its modification or deletion. The Fact Layer enforces the single-assignment constraint architecturally: the write interface exposes no \texttt{UpdateParentPointer} or \texttt{DeleteParentPointer} method.
    \item \textbf{Forensic Ledger Entry:} The Fact Layer appends an entry to the forensic ledger containing:
    \begin{itemize}
        \item The child identifier $c$ and parent identifier $p$.
        \item The depth $d_c$.
        \item A timestamp $\tau_{\text{spawn}}$ from the Purpose Layer's atomic clock.
        \item A cryptographic hash of the mandate $H(\text{mandate}_c)$.
        \item The chain hash $H_{\text{prev}}$ linking this entry to the previous ledger entry, forming a cryptographically linked chain analogous to a Merkle tree or blockchain-style log.
    \end{itemize}
    \item \textbf{Child Worksheet Sealing:} The child $c$ is initialized with a sealed Worksheet containing:
    \begin{itemize}
        \item The complete parent pointer $ParentPointer(p, c)$, sealed at initialization and thereafter unreadable from the application layer.
        \item The child's declared mandate and convergence criteria.
        \item A sealed reference to the parent's Purpose Layer for upstream accountability escalation.
    \end{itemize}
\end{enumerate}

The immutability of the parent pointer created in this phase is the central guarantee of RSP. It is not enforced by access control or coding convention — it is architecturally enforced by the Fact Layer's append-only write interface. Any attempt to alter the parent pointer at the application layer, middleware layer, or database layer is rejected by the Fact Layer without exception or appeal.

\bigskip
\noindent\textbf{Phase 4: Child Activation with Immutable Parent Reference.}

After the Fact Layer confirms the successful atomic write, the child $c$ is activated. The child $c$ is initialized with its immutable parent reference and operates within the following structural guarantees:

\begin{enumerate}
    \item \textbf{Upstream Accountability:} If $c$ encounters a state it cannot resolve within its bounded mandate, $c$ escalates to its parent $p$ via the Purpose Layer's upstream accountability mechanism. If $p$ is itself drifted (Section~\ref{sec:autonomous-drift}), the escalation chain continues upward until it reaches a non-drifted anchor. The drift of $p$ does not sever $c$'s accountability chain; it triggers an escalation that propagates toward the human anchor.
    \item \textbf{No Re-parenting:} Once activated, $c$ cannot re-parent itself, cannot spawn with a different parent, and cannot sever its parent pointer under any condition. The Fact Layer enforces this immutability structurally.
    \item \textbf{Chain Termination at Human:} The parent chain of $c$, traced recursively through $ParentPointer$, is guaranteed to terminate at the root human $h_{\text{root}}$, who holds the null parent sentinel $ParentPointer(\bot, h_{\text{root}})$. This termination is enforced at the Fact Layer for every spawn event and is verifiable at any time by any party with read access to the forensic ledger.
\end{enumerate}

The activation of $c$ is not contingent on runtime conditions, network availability, or parent continued operation. $c$ holds a structurally immutable reference to $p$, which holds a structurally immutable reference to its own parent, and so on. The chain is indestructible by design.

% -------------------------------------------------------
\subsection{Depth Management}

Depth management in RSP serves two distinct functions that are often conflated in practice: \emph{resource governance} and \emph{drift containment}. RSP separates these concerns architecturally.

\bigskip
\noindent\textbf{Maximum Spawn Depth ($\lambda_{\text{max\_depth}}$).}

The system enforces a hard upper bound on spawn depth via the Bounds Engine. This parameter is configurable at system initialization and may be adjusted by a human administrator, but the adjustment itself requires Purpose Layer authorization and is logged to the forensic ledger.

Formally, for any spawn event creating child $c$ with parent $p$:

\begin{equation}
    \text{DepthCheck}(c) \equiv 
    \begin{cases}
    \text{Approved} & \text{if } d_c \leq \lambda_{\text{max\_depth}} \\
    \text{Refused} & \text{if } d_c > \lambda_{\text{max\_depth}}
    \end{cases}
\end{equation}

where $d_c = \text{Depth}(p) + 1$ and $\text{Depth}(h_{\text{root}}) = 0$.

The key insight from the drift analysis (Section~\ref{sec:drift-problem}) is that $\lambda_{\text{max\_depth}}$ functions as a \textit{blast radius limiter}, not a drift preventer. A depth limit does not prevent drift from occurring at any depth below the cap — it only limits how many generations of drifted agents can accumulate before the system considers the chain exhausted. Drift containment requires that every node in the chain be non-drifted, which is enforced by the Purpose Layer's mandate coherence check at each spawn event, not by depth limits alone.

\bigskip
\noindent\textbf{Spawn Refusal at Depth Limit.}

When a spawn event is refused at the depth limit, the following occurs:

\begin{enumerate}
    \item The Bounds Engine returns \texttt{SpawnRefusal} with the reason code \texttt{DEPTH\_LIMIT\_EXCEEDED}.
    \item The Purpose Layer logs the refused event to the forensic ledger as a denied authorization, including the proposed child ID, parent ID, proposed depth, and timestamp.
    \item The parent $p$ receives a \texttt{SpawnRefusal} response, enabling it to pursue alternative strategies: decomposing the task differently, escalating to its own parent for a broader authorization, or returning a \texttt{TASK\_INCOMPLETE} signal to the human anchor.
    \item No partial spawn record is created. The refused child $c$ does not exist in the system.
\end{enumerate}

This refusal protocol ensures that depth limit enforcement is auditable — every refusal is recorded — and that parent agents have a defined response path when encountering the limit, rather than failing silently or spawning unauthorized children.

\bigskip
\noindent\textbf{Convergence Criteria.}

A convergence criterion is a formally specified condition under which a child agent terminates and returns its outputs to its parent. It is not a soft guideline — it is a structural requirement enforced by the Bounds Engine before any spawn is authorized. Without a convergence criterion, a spawned child has no defined exit condition, and its parent cannot determine whether the child's task is complete, stalled, or failed.

Convergence criteria in RSP take three forms:

\begin{itemize}
    \item \textbf{Result Delivery:} The child delivers a defined output artifact (data structure, document, decision record) to the parent, satisfying the task mandate. Upon delivery, the child enters a suspended state pending parent acknowledgment.
    \item \textbf{Exception Signal:} The child encounters a state that falls outside its mandate and cannot be resolved by escalation (e.g., a logical impossibility, a resource exhaustion with no escalation path). This is treated as a convergence event with an exception output, not as a failure.
    \item \textbf{Timeout:} The child exceeds a configured maximum active duration $\tau_{\text{max}}$ without delivering a result. Timeout is a convergence event with a timeout output — it does not mean the child was malfunctioning; it means the task framing was incorrect or the resources provided were insufficient.
\end{itemize}

The enforcement of convergence criteria at spawn time is what distinguishes RSP's depth management from the depth limits used in conventional agent frameworks. Depth limits in conventional systems are a blunt instrument: they cap the number of spawn levels but do not ensure that any given level has a defined exit. RSP's convergence criteria ensure that every spawned child has a defined termination condition \textit{before it is created}, converting the spawn event from an open-ended fork into a bounded delegation.

% -------------------------------------------------------
\subsection{Chain Termination Verification}

RSP requires that every agent, at any depth, have a verifiable chain terminating at a human principal. This property — chain termination at human — is established at spawn time, maintained through the agent's lifetime, and verifiable at any point via the forensic ledger.

\bigskip
\noindent\textbf{Formal Chain Termination Property.}

For any agent $a$ in the system, define the \emph{accountability chain} $\text{Chain}(a)$ recursively as:

\begin{equation}
    \text{Chain}(a) \equiv 
    \begin{cases}
        \emptyset & \text{if } a = h_{\text{root}} \\
        \{ ParentPointer(\text{parent}(a), a) \} \cup \text{Chain}(\text{parent}(a)) & \text{if } a \neq h_{\text{root}}
    \end{cases}
\end{equation}

The RSP chain termination guarantee is:

\begin{equation}
    \forall a:\; \text{Chain}(a) \neq \emptyset \implies h_{\text{root}} \in \text{Chain}(a)
\end{equation}

That is, every agent's chain contains the root human. This is enforced by the Fact Layer at spawn time: every \texttt{RecordSpawnEvent} operation must satisfy this property, and the operation is rejected if it would create a chain that does not terminate at a human.

\bigskip
\noindent\textbf{Verification Protocol.}

Any party — human, auditor, external system — can verify the chain termination property for any agent $a$ using the following protocol:

\begin{enumerate}
    \item Retrieve $a$'s current parent pointer from the Fact Layer's immutable registry.
    \item Recursively retrieve each parent's parent pointer, yielding the ordered set $\text{Chain}(a)$.
    \item Verify that the terminal element of the chain is $h_{\text{root}}$ with parent sentinel $\bot$.
    \item Verify the cryptographic integrity of each ledger entry using the chain hash $H_{\text{prev}}$.
    \item Verify that each entry's mandate hash $H(\text{mandate}_i)$ is consistent with the declared task at each depth.
\end{enumerate}

This verification protocol is computationally lightweight — it requires $O(d)$ ledger reads where $d$ is the depth of $a$ — and requires no access to the internal state of any agent in the chain. The forensic ledger is a public read surface: the ledger itself is append-only and cryptographically sealed, so any reader can verify chain integrity without trusting the agents in the chain.

The practical significance is that chain termination verification can be performed by a human supervisor in seconds, without specialized tooling or access to agent runtime state. This is what makes RSP auditable in practice, not merely in theory: the audit is a simple, deterministic read operation, not a complex forensic reconstruction requiring access to proprietary agent logs.

% -------------------------------------------------------
\subsection{Protocol Properties}

The Recursive Spawning Protocol satisfies the following formal properties:

\bigskip
\noindent\textbf{Single-Assignment Parent Pointer.} The Fact Layer's $ParentPointer(p, c)$ relation is assigned exactly once per spawn event and is thereafter permanently immutable. No operation in the protocol, no runtime condition, and no administrative action can modify or delete a parent pointer after assignment.

\bigskip
\noindent\textbf{Atomic Spawn.} Every spawn event is atomic across the three layers: all four phases succeed together, or the spawn event is rejected in its entirety. There is no partially recorded spawn event in the forensic ledger. This is the equivalent of a database transaction with serializable isolation applied to spawn events.

\bigskip
\noindent\textbf{Drift-Non-Amplifying Spawn.} Because the Purpose Layer validates mandate coherence and the Bounds Engine enforces convergence criteria before any child is created, RSP prevents the creation of children that have no defined exit or that carry a mutated mandate. This breaks the drift amplification condition: a drifted parent cannot spawn a child that propagates drift, because the drifted condition itself prevents the Purpose Layer's mandate coherence check from passing.

\bigskip
\noindent\textbf{Forensic Ledger Completeness.} Every spawn event — authorized or refused — is recorded in the forensic ledger. The ledger is append-only and cryptographically linked, making it tamper-evident. The completeness guarantee means that any audit of the system can reconstruct the full spawn history from genesis to current state, with cryptographic proof of integrity.

\bigskip
\noindent\textbf{Human-Anchored Termination.} Every agent's accountability chain, traced via the immutable parent pointers, terminates at the root human $h_{\text{root}}$. This is not a probabilistic property — it is a deterministic guarantee enforced by the Fact Layer at every spawn event. There exists no agent in the system whose chain does not terminate at a human.

% -------------------------------------------------------
\subsection{Operational Implications}

RSP's guarantees have direct operational consequences for multi-agent system design.

First, the depth of any task decomposition is bounded by $\lambda_{\text{max\_depth}}$, but the more important bound is the convergence criterion requirement: each child must have a defined exit condition before it is created. This means that system designers must frame tasks as having specified termination conditions — a discipline that itself reduces the risk of unbounded agent proliferation.

Second, the forensic ledger provides a complete, immutable history of every spawn event in the system. This history is not for post-hoc reconstruction of what happened — it is for real-time verification that the system is operating within its authorized structure. A human supervisor can at any time query any agent's chain and verify that the authorization path to a human is intact.

Third, the separation of Purpose Layer authorization, Bounds Engine depth/convergence checking, and Fact Layer recording means that each layer can be audited, upgraded, or replaced independently without invalidating the guarantees of the other layers. The immutable parent pointer is the stability anchor: the Fact Layer's guarantees hold regardless of what Purpose Layer policy is in force or what Bounds Engine configuration is applied.

Fourth, the atomic spawn property means that system failures (network loss, process crash, power interruption) cannot produce partial spawn records. If a spawn event is interrupted between phases, the Fact Layer's write is not committed and no record exists — the spawn is as if it never happened. This prevents the creation of \emph{orphan agents} that exist in the system without a complete authorization record.

RSP transforms recursive spawning from a provenance gap — the most dangerous failure mode in hierarchical multi-agent systems — into the most precisely audited and structurally guaranteed operation in the system. Every spawn is a first-class accountability event, recorded immutably, with a verifiable chain to a human principal.